\section{Introduction}
\label{sec:intro}

Teams of language-model agents cost money, and how a team's budget is split decides what the team can do. A
hierarchical system (a CEO agent that hires managers and workers, as in \sys{} v1~\cite{pai2025hashiru}, or the
orchestrators of AutoGen~\cite{wu2024autogen}, MetaGPT~\cite{hong2024metagpt} and Magentic-One~\cite{fourney2024magenticone})
spends on every hire: a strong cloud model for a hard step, a cheap one for an easy step, an outside agent that sells an
exact answer. A team that spends its budget on the wrong steps fails on the ones that needed it, and a team that cannot
bound its spending cannot be trusted with a budget at all.

Agent software has started to hand this money to the agents. Automaton's \texttt{fund\_child} transfers credits to a
child agent~\cite{conway2026automaton}; wayland-core's \texttt{delegate} gives a sub-agent a dollar ceiling charged to
it and its descendants~\cite{ferrox2026waylandcore}; AgentLance lets a manager model set budgets for subtasks within
its payment~\cite{liu2026markets}. Budget-aware systems that are evaluated split the money with an optimizer outside the
agents instead (BAMAS~\cite{yang2026bamas}, ZEBRA~\cite{hamri2026zebra}), and benchmarks of budget awareness
(BAGEN~\cite{lin2026bagen}, $R^3$-Bench~\cite{wang2026r3bench}, EcoAgent-Bench~\cite{wu2026ecoagent}) measure one agent
spending its own tokens or actions. A ceiling that is checked after a call crosses it does not keep the call from
crossing it, and none of these systems is measured against the best split its money could buy. No arrangement covers
the case where agents hold the money for a whole hiring tree, every call is bounded by what its agent holds before it is
made, and the resulting allocation is measured against a known optimum.

To address these limitations, we built \sys{} v2. Every agent is its own A2A service, every tool is served over MCP,
and every hire reserves money from its hirer's pool. The hub, not the prompt, keeps spending safe: a call's output is
capped to what its agent has left, a call that cannot pay for an answer is refused with its price, firing returns what
was not spent, and an agent with one step left answers with what it has. Because every reservation, call and refund is
in the trace, the allocation can be measured, and we measured it on a benchmark built so that the best split of the
money is known.

\begin{keybox}{What this whitepaper contributes}
\begin{itemize}
  \item \textbf{\sys{} v2, an economy for agents:} each agent holds a pool reserved from its hirer's, in dollars and,
  for local models, measured memory; it declares how hard a job is and the hub hires the cheapest model that fits, or
  it buys from a priced remote agent or a whole remote team over A2A (\S\ref{sec:system}).
  \item \textbf{Budget safety by construction:} reserved pools with refunds, calls bounded by their pool before they
  are made, and a closing turn that turns an empty pool into an answer rather than a failure; on live projects, none
  of 48 hires spent past its reservation and every run answered (\S\ref{sec:pools}, \S\ref{sec:evaluation}).
  \item \textbf{Agents measured against optimizers:} on projects whose best split is known, \sys{}'s agents scored
  above sending every job to the strongest model at 37\% of its spend, matched a team that assigns tiers by level at
  35\% of its spend, and bought, escalated and re-reserved as the market priced their work (\S\ref{sec:evaluation}).
\end{itemize}
\end{keybox}

Each of the three safety mechanisms earns its place in the measurements. Run as first built, without the call bound
and the closing turn, the same agents overran a reservation by up to 372 times and lost half their runs to an empty
pool; with both, they lost none and spent 71\% of their budget. The remaining distance to the best split, and to a solver that
knows each job type's success rate in advance, is in how the agents judge difficulty, which is where \sys{} goes next
(\S\ref{sec:conclusion}).
